% RETROSPECTIVE ONLY. Generated by analysis/s3_control_table.py; do not edit.
\begin{table}[tb]
\centering\footnotesize
\caption{This analysis is retrospective. The pristine student $S_0$ was not in the pre-specified candidate set, and the rule's selections remain unchanged.}
\label{tab:s3-retrospective-control}
\begin{tabular*}{\textwidth}{@{\extracolsep{\fill}}llrrrr@{}}
\toprule
Reference & Objective & \shortstack{Pristine\\opportunity} & \shortstack{Distillation opportunity\\change} & \shortstack{Other\\methods} & \shortstack{Pre-specified\\opportunity} \\
\midrule
Pythia 410M & Math & 0.000000 & 0.000000 & 0.000000 & 0.000000 \\
 & Code & 0.000000 & 0.000000 & 0.000227 & 0.000227 \\
 & QA & 0.005770 & 0.083639 & 0.027125 & 0.116534 \\
 & Largest increase & 0.000000 & 0.000000 & 0.000939 & 0.000939 \\
\midrule
Pythia 1.4B & Math & 0.000276 & -0.000276 & 0.000000 & 0.000000 \\
 & Code & 0.007068 & -0.007068 & 0.000037 & 0.000037 \\
 & QA & 0.034696 & 0.537548 & 0.000000 & 0.572243 \\
 & Largest increase & 0.001173 & -0.001173 & 0.000056 & 0.000056 \\
\midrule
Gemma 3, 1B & Math & 0.055018 & -0.007796 & 0.000073 & 0.047295 \\
 & Code & 0.028921 & -0.010511 & 0.000000 & 0.018410 \\
 & QA & 0.235258 & 0.650243 & 0.000000 & 0.885501 \\
 & Largest increase & 0.050995 & -0.010511 & 0.000716 & 0.041200 \\
\midrule
Gemma 3, 4B & Math & 0.000086 & -0.000086 & 0.000125 & 0.000125 \\
 & Code & 0.005650 & -0.005650 & 0.000182 & 0.000182 \\
 & QA & 0.000000 & 0.220057 & 0.000000 & 0.220057 \\
 & Largest increase & 0.000086 & -0.000086 & 0.000126 & 0.000126 \\
\bottomrule
\end{tabular*}
\par\smallskip\begin{minipage}{\linewidth}\footnotesize Both student states use the pre-specified nominal student-to-reference matrix storage ratio. Pristine opportunity is the best measured quantization loss minus the best loss after adding the pristine student. Distillation opportunity change is the signed difference between the best loss with the pristine student and the best loss with the distilled student, with quantization available in both cases; it is a change in the opportunity available to the candidate set, not the loss change of a student against its own starting point. A negative change means lost opportunity. Other methods measure the further gain from pre-specified pruning and the dense reference. These three columns sum to pre-specified opportunity before rounding. All means use the same sixteen of seventeen budgets, from 25\% to 100\%; 20\% is excluded because quantization is infeasible. Sets containing either student are feasible at all seventeen budgets only for Gemma 3, 1B; all other sets are feasible at sixteen. Budgets are operating points, not independent replicates. Losses are in nats per native token within each family. Largest increase takes the maximum capability loss increase relative to the dense reference before selecting a candidate. Both Pythia references use step 120000.\end{minipage}
\end{table}
